% Part: incompleteness
% Chapter: theories-computability
% Section: introduction

\documentclass[../../../include/open-logic-section]{subfiles}

\begin{document}

\olfileid{inc}{tcp}{int}

\olsection{Pambuka}

\begin{editorial}
Bagean iki perlu ditulis maneh.
\end{editorial}

Kita wis nduweni:
\begin{enumerate}
\item definisi teges fungsi kang representabel ing
  $\Th{Q}$ (\olref[req][int]{defn:representable-fn})
\item definisi teges relasi kang representabel ing
  $\Th{Q}$ (\olref[req][rel]{defn:representing-relations})
\item teorema kang nyatakake yen fungsi representabel ing $\Th{Q}$
  persis fungsi komputabel
  (\olref[req][int]{thm:representable-iff-comp})
\item teorema kang nyatakake yen relasi representabel ing $\Th{Q}$
  persis relasi komputabel
  (\olref[req][rel]{thm:representing-rels})
\end{enumerate}

Sawijining {\em teori} iku himpunan ukara kang katutup tumrap deduksi,
yaiku nduweni sipat: saben $T$ mbuktekake $!A$, mula $!A$ kalebu ing
$T$. Teori bisa dianggep minangka kumpulan ukara bebarengan karo
kabeh perkara kang diakibatake dening ukara mau. Wiwit saiki, kita
nggunakake $\Th{Q}$ kanggo nyebut {\em teori} kang dumadi saka
himpunan ukara kang bisa diturunake saka wolung aksioma ing
\olref[req][int]{sec}. Elinga yen kita bisa ngodeake formula
ing basa $\Th{Q}$ minangka wilangan; yen $!A$ formula kaya mangkono,
tetepna $\Gn{!A}$ minangka wilangan kang ngodeake $!A$.
Kanthi pangodean iki, saiki kita bisa takon apa maneka himpunan
formula komputabel utawa ora.

\end{document}
